\answersfor{ch:trainloop}

\begin{solution}{ex:trainloop:1}
$L'(w) = 2(w - 3)$, $\eta = 0.25$, $w_0 = 0$.
\begin{align*}
\text{Step 1: } g &= 2(0 - 3) = -6, & w_1 &= 0 - 0.25 \times (-6) = 1.5\\
\text{Step 2: } g &= 2(1.5 - 3) = -3, & w_2 &= 1.5 - 0.25 \times (-3) = 2.25\\
\text{Step 3: } g &= 2(2.25 - 3) = -1.5, & w_3 &= 2.25 - 0.25 \times (-1.5) = 2.625
\end{align*}
The distances to 3 are $-3, -1.5, -0.75, -0.375$: each is multiplied by $1 - 2\eta = 1 - 0.5 =
0.5$. One step lands exactly on 3 when the factor is 0, that is $1 - 2\eta = 0$, $\eta = 0.5$:
from $w_0 = 0$, $w_1 = 0 - 0.5 \times (-6) = 3$. (For this loss the curvature is $L'' = 2$, and
the one-step learning rate is $1/L''$.)
\end{solution}

\begin{solution}{ex:trainloop:2}
$g = 2$, $g^2 = 4$, $\beta_1 = 0.9$, $\beta_2 = 0.999$.

\textbf{$t = 1$.}
\begin{align*}
m_1 &= 0.9 \times 0 + 0.1 \times 2 = 0.2\\
v_1 &= 0.999 \times 0 + 0.001 \times 4 = 0.004\\
\hat{m}_1 &= 0.2 / 0.1 = 2\\
\hat{v}_1 &= 0.004 / 0.001 = 4, \quad \sqrt{4} = 2\\
\text{step} &= 10^{-3} \times 2/2 = 10^{-3}
\end{align*}
\textbf{$t = 2$.}
\begin{align*}
m_2 &= 0.9 \times 0.2 + 0.1 \times 2 = 0.18 + 0.2 = 0.38\\
v_2 &= 0.999 \times 0.004 + 0.001 \times 4 = 0.003996 + 0.004 = 0.007996\\
\hat{m}_2 &= 0.38 / 0.19 = 2\\
\hat{v}_2 &= 0.007996 / 0.001999 = 4, \quad \sqrt{4} = 2\\
\text{step} &= 10^{-3} \times 2/2 = 10^{-3}
\end{align*}
Bias correction recovers the true mean (2) and mean square (4) exactly for a constant gradient,
and the step is $\eta$.

\textbf{Without correction at $t = 1$:} $m_1/\sqrt{v_1} = 0.2/\sqrt{0.004} = 0.2/0.063246 =
3.1623$, so the step would be $3.16 \times 10^{-3}$, more than three times too large
($\sqrt{1000}/10 = 31.62/10 = 3.162$).
\end{solution}

\begin{solution}{ex:trainloop:3}
$\eta(t) = 10^{-6} + 0.0004995\,(1 + \cos(\pi t/200))$.

\textbf{$t = 25$:} $\pi \times 25/200 = \pi/8$, $\cos(\pi/8) = 0.923880$.
\begin{align*}
1 + 0.923880 &= 1.923880\\
0.0004995 \times 1.923880 &= 0.00096098\\
0.000001 + 0.00096098 &= 0.00096198 = 9.620 \times 10^{-4}
\end{align*}
\textbf{$t = 175$:} $\pi \times 175/200 = 7\pi/8$, $\cos(7\pi/8) = -0.923880$.
\begin{align*}
1 - 0.923880 &= 0.076120\\
0.0004995 \times 0.076120 &= 0.000038022\\
0.000001 + 0.000038022 &= 0.000039022 = 3.902 \times 10^{-5}
\end{align*}
\textbf{Rate of change.} $d\eta/dt = -\frac{\pi}{400}(0.000999)\sin(\pi t/200)$. Since $\sin(\pi/8) =
\sin(7\pi/8) = 0.382683$, the rate is the same at both epochs: $-0.0078540 \times 0.000999
\times 0.382683 = -3.00 \times 10^{-6}$ per epoch. The cosine curve is symmetric about $t = 100$.
In \emph{relative} terms, though, the change at $t = 175$ is far larger: $3.0 \times 10^{-6}$ is
0.3\% of $9.6 \times 10^{-4}$ but 7.7\% of $3.9 \times 10^{-5}$.
\end{solution}

\begin{solution}{ex:trainloop:4}
$\norm{(1, 2, 2)} = \sqrt{1 + 4 + 4} = \sqrt{9} = 3$.

\textbf{To 1.5:} $\min(1, 1.5/3) = 0.5$, so $\vect{g} \to (0.5, 1, 1)$, whose length is
$\sqrt{0.25 + 1 + 1} = 1.5$.

\textbf{To 5:} $\min(1, 5/3) = 1$: unchanged, $(1, 2, 2)$.

\textbf{Would clipping at 1.0 change this repository's training?} Adam divides each parameter's
gradient by a running estimate of its size, so if every step's gradient were clipped by the same
factor, $\hat{m}/\sqrt{\hat{v}}$ would be unchanged and so would training. Clipping changes
things only on the steps where the gradient norm is unusually large: it shrinks those gradients
before they enter $m$ and $v$, so one wild batch moves the parameters and inflates $v$ less.
Whether a threshold of 1.0 is ever active depends on the typical whole-model gradient norm,
which the repository does not log; you would measure it before choosing a threshold.
\end{solution}

\begin{solution}{ex:trainloop:5}
\textbf{Rows} (height 400):
\begin{align*}
0 + 128 &= 128 \le 400 \text{ keep } 0;\quad 64 + 128 = 192 \text{ keep } 64;\quad 128 + 128 = 256 \text{ keep } 128;\\
192 + 128 &= 320 \text{ keep } 192;\quad 256 + 128 = 384 \le 400 \text{ keep } 256;\\
320 + 128 &= 448 > 400 \text{: snap to } 400 - 128 = 272, \text{ stop.}
\end{align*}
Rows $[0, 64, 128, 192, 256, 272]$: 6.

\textbf{Columns} (width 300): as in \cref{we:trainloop:plan}, $[0, 64, 128, 172]$: 4.

\textbf{Patches:} every row with every column, $6 \times 4 = 24$. The real generator returns the
same \tagC.

\textbf{A duplicate:} any height where $(\text{height} - 128)$ is a multiple of 64, for example
192: rows 0 and 64 are kept ($64 + 128 = 192 \le 192$), then row 128 gives $256 > 192$ and snaps
to $192 - 128 = 64$, already present. Heights 128, 256 and 320 do the same.
\end{solution}

\begin{solution}{ex:trainloop:6}
Near 0.2 a bfloat16 number has 8 significant bits, so neighbouring representable values are
$2^{-3} \times 2^{-7} = 2^{-10} = 0.000977$ apart: 0.2001953, then 0.2011719. Both 0.2 and
0.2004 round to 0.2001953, so their difference, 0.0004, becomes exactly 0. A good model's errors
are of this size, so a loss computed in bfloat16 would be blind to precisely the errors training
is trying to remove.

The SpectralEarth trainer runs only the model's forward pass under bfloat16 autocast and
converts the output to float32 (\code{out.float()}) before the loss. The truth is never rounded,
and the subtraction, absolute value, angle and averaging run in float32. (The prediction itself
still carries the rounding of the network's last bfloat16 operations, in the network's
normalised units.)
\end{solution}
